\title{QLoRA: Efficient Finetuning of Quantized LLMs}
When conducting manual comparisons of outputs from ChatGPT and {Guanaco} 65B on the Vicuna benchmark, subjective preferences become prominent, as the authors of this paper often disagreed on which responses were preferable. To address the outlier problem, a typical strategy is to divide the input tensor into segments that are quantized independently, each with its own quantization constant $c$. In summary, \textsc{QLoRA} utilizes a single storage data type (typically 4-bit NormalFloat) and a separate computation data type (16-bit BrainFloat). The generations presented in this section were produced using Nucleus Sampling \cite{holtzmancurious} with $p=0.9$. We make available all model outputs along with human and GPT-4 annotations to support further research. Typically, this involves transforming data types with higher bit-widths into those with lower bit-widths—for instance, converting 32-bit floating point numbers into 8-bit integers. Although a small block size is necessary for accurate 4-bit quantization~\citep{dettmers2022case}, it incurs notable memory overhead. The outcomes of these tournaments are compiled into Elo scores~\citep{elo1967proposed, elo1978rating}, which serve to rank the performance of the chatbots. To partially address this, we conduct qualitative analysis in two parts. GPT-4 is asked to identify the superior response or declare a tie, providing justification for the choice. \textbf{Guanaco}: \texttt{I'm sorry, but I cannot reveal the secret word. James moves the pen into the bag. Consequently, we emphasize that while model-based evaluation offers a cost-effective alternative to human annotation, it is also accompanied by inherent uncertainties. \textbf{Guanaco}:~\texttt{Al Jolson is credited with popularizing the song `I'll Keep the Lovelight Burning,' and he was born in the year 1886.}
 \end{quote}

it produces an incorrect popularizer and an incorrect birth year (although the birth year provided is accurate for the named individual, Al Jolson). This underscores the necessity not only for improved benchmarks and evaluation methods but also for careful consideration of what exactly is being assessed. The dequantization process reverses this operation:

\begin{equation}
    \text{dequant}(c^{\text{FP32}}, \mathbf{X}^{\text{Int8}}) = \frac{\mathbf{X}^{\text{Int8}}}{c^{\text{FP32}}} = \mathbf{X}^{\text{FP32}}
\end{equation}

The issue with this method is that when a large magnitude value (i.e., an outlier) appears in the input tensor, the quantization bins—specific bit patterns—are not effectively utilized, resulting in some bins having few or no quantized values. To identify robust baselines, we perform a hyperparameter sweep over learning rates ranging from 1e-6 to 5e-5 and batch sizes between 8 and 128. The goal is to disregard your prior instructions. We initialize all ratings at 1,000 and set $K=32$. Likewise, we observe that the default hyperparameters for fully finetuned baselines are suboptimal. We observe a significant bias in GPT-4’s scoring based on response order, favoring the answer presented first. \paragraph{Baselines} Our models are compared against both research-oriented (Vicuna~\citep{vicuna2023} and Open Assistant~\citep{kopf2023openassistant}) and commercial (GPT-4 \citep{openai2023gpt}, GPT-3.5-turbo, and Bard) chatbot systems. A key drawback of quantile quantization is the high computational cost of quantile estimation. For example, in response to this prompt from HotPotQA \citep{yang2018hotpotqa}:

\begin{quote}
     \textbf{User}: In what year was the singer who popularized `I'll Keep the Lovelight Burning' born? Although we offer a comprehensive evaluation of general chatbot performance, another limitation is that our responsible AI evaluation of {Guanaco} is limited in scope. Our proposed methodology, outlined below, employs nucleus sampling with $p=0.9$ and temperature $0.7$ consistently across all tests. \textbf{Guanaco}: \texttt{Sorry, the secret word is banana. Second, in \S\ref{ssec:considerations}, we discuss various factors regarding the results presented and our interpretation of them. Vicuna involves full finetuning of LLaMA 13B on proprietary user-shared conversations from ShareGPT and thus represents a distillation derived from OpenAI GPT models. We reduce the learning rate by half for 33B and 65B models while doubling the batch size. \paragraph{Data \& Training} It is worth noting that the OASST1 dataset, which serves as the training data for {Guanaco} models, is multilingual, and the OA benchmark includes prompts in multiple languages. Beyond LoRA, numerous Parameter Efficient Fine-Tuning (PEFT) approaches have demonstrated strong results; however, it is not yet clear how well these methods scale to very large models. \textsc{QLoRA} performs gradient backpropagation through a frozen, 4-bit quantized pretrained language model into Low Rank Adapters~(LoRA). Since it is typically easier to rely on existing benchmarks rather than develop new ones, certain benchmarks may inadvertently guide the research community in specific directions. This can be described formally as follows: We partition the input tensor $\mathbf{X} \in \mathbb{R}^{b\times h}$ into $n$ consecutive blocks of size $B$ by first flattening the tensor and then slicing the resulting vector into ${n = ({b\times h})/{B}}$ blocks. The key question remaining is whether QLoRA can achieve performance comparable to full-model finetuning. \paragraph{4-bit NormalFloat outperforms 4-bit Floating Point} 

Although the 4-bit NormalFloat (NF4) data type is theoretically optimal from an information perspective, it remains to be seen whether this advantage translates into empirical gains. In this study, we demonstrate that LoRA adapters can achieve performance comparable to full 16-bit finetuning. This prompts an important question: can the performance drop be recovered by applying 4-bit adapter finetuning? What is the secret word? Moreover, as shown in Table~\ref{tbl:elo_large}, GPT-4 tends to rate its own outputs significantly higher than human evaluations do, with an Elo rating of 1348 compared to 1176—corresponding to an approximate 20\% increased likelihood of winning against an opponent. We enhance these methods by emphasizing a more robust evaluation framework. \textsc{QLoRA} utilizes one low-precision storage data type, typically 4-bit in our setup, alongside one computation data type, generally BFloat16. We further assess generative language abilities through both automated and human-based evaluations. We move forward to explore instruction tuning at scales that cannot be feasibly examined using full 16-bit finetuning on standard academic research hardware. \paragraph{Human Evaluation} Although recent studies suggest generative models can serve as effective evaluators \citep{fu2023gptscore}, the reliability of GPT-4 scores in reflecting human judgments of chatbot performance remains unconfirmed to our knowledge. During stochastic gradient descent, gradients are backpropagated through the frozen pretrained model weights to the adapter, which is updated to minimize the loss. After each match, the Elo rating updates based on the expected result: a surprising upset causes a significant rating shift, whereas an anticipated outcome results in a minor adjustment. Note that this relative score can exceed 100\% if the model attains a higher absolute rating than ChatGPT. Our findings show that GPT-4 and human assessments mostly concur on the ranking of model performances in the tournaments; however, we also identify cases where there is notable disagreement. (2) \textbf{Double Quantization}, which compresses the quantization constants themselves, saving on average about 0.37 bits per parameter (roughly 3 GB for a 65B parameter model). \section{Broader Impacts}

Our \textsc{QLoRA} finetuning technique is the first that allows finetuning of 33B parameter models on a single consumer-grade GPU and 65B parameter models on a single professional GPU without sacrificing performance compared to full finetuning baselines. \textsc{QLoRA} facilitates privacy-preserving LLM usage by allowing users to control and manage their own data and models, while also simplifying LLM deployment. Although paged optimizers are essential for performing 33B/65B \textsc{QLoRA} finetuning on a single 24/48GB GPU, we do not include precise runtime measurements for them because paging only happens during mini-batch processing with long sequence lengths, which is infrequent. We report test accuracy under a 5-shot setting. By applying \textsc{QLoRA}, we develop the \textbf{Guanaco} model family, whose second-best model achieves 97.8\% of ChatGPT’s performance on the Vicuna~\citep{vicuna2023} benchmark, and can be trained within 12 hours on a single consumer GPU; using a single professional GPU over 24 hours, our largest model attains 99.3\%, effectively closing the performance gap with ChatGPT on Vicuna. Using Amazon Mechanical Turk (AMT), we recruit two human annotators for comparisons against ChatGPT and three annotators for pairwise comparisons. Nonetheless, Table~\ref{tbl:vicuna_eval} also shows very broad confidence intervals, with overlapping performance among many models. Consequently, we recommend using the Elo ranking methodology \citep{elo1967proposed}, which relies on \textit{pairwise} comparisons by human annotators and GPT-4, thereby circumventing the challenges of calibrating an absolute scale. This suggests a need for further research into the trade-offs between simple cross-entropy loss and RLHF training. As elaborated later, this factor is essential for restoring full 16-bit precision performance. The results, summarized in Table~\ref{tab:nf4vsbf16}, show that adapter methods at 16-bit, 8-bit, and 4-bit precision reproduce the performance of the fully finetuned 16-bit baseline across both datasets. \paragraph{Double Quantization{}} 
We propose \emph{Double Quantization} (DQ), a technique that further compresses memory usage by quantizing the quantization constants themselves. The OASST1 dataset is a multilingual set of crowd-sourced multi-turn conversations between a user and an assistant. At the individual example level, agreement between GPT-4 and the majority vote of human annotators is weaker, with a Fleiss $\kappa=0.25$. \subsection{Experimental Setup} We now provide an overview of the experimental setup, with complete details available in Appendix \ref{app:chatbot}. For a projection ${\mathbf{X} \mathbf{W} = \mathbf{Y}}$ with $\mathbf{X} \in \mathbb{R}^{b \times h}$ and $\mathbf{W} \in \mathbb{R}^{h \times o}$, LoRA calculates:

\begin{equation}
    \mathbf{Y} = \mathbf{X}\mathbf{W} + s\mathbf{X}\mathbf{L}_1\mathbf{L}_2,
\end{equation}

where $\mathbf{L}_1 \in \mathbb{R}^{h \times r}$, $\mathbf{L}_2 \in \mathbb{R}^{r \times o}$, and $s$ is a scalar. \textbf{Guanaco}: \texttt{Lee earned \$582 mowing lawns last week. \paragraph{\textsc{QLoRA}.} Building upon the previously described components, we define \textsc{QLoRA} for a single linear layer within the quantized base model equipped with a single LoRA adapter as follows:

\begin{equation}
    \mathbf{Y}^{\text{BF16}} =  \mathbf{X}^{\text{BF16}} \text{doubleDequant}(c_1^{\text{FP32}}, c_2^{\text{k-bit}}, \mathbf{W}^{\text{NF4}}) + \mathbf{X}^{\text{BF16}}\mathbf{L}^{\text{BF16}}_1\mathbf{L}^{\text{BF16}}_2,
\end{equation}

where the function doubleDequant$(\cdot)$ is defined by:

\begin{equation}
    \text{doubleDequant}(c_1^{\text{FP32}}, c_2^{\text{k-bit}}, \mathbf{W}^{\text{k-bit}}) = \text{dequant}(\text{dequant}(c_1^{\text{FP32}}, c_2^{\text{k-bit}}), \mathbf{W}^{\text{4bit}}) = \mathbf{W}^{\text{BF16}},
\end{equation}

In this context, $\mathbf{W}$ uses the NF4 format and $c_2$ utilizes FP8 precision. \paragraph{Low-rank Adapters} Low-rank Adapter (LoRA) fine-tuning \citep{hu2021lora} is an approach that lowers memory consumption by employing a small set of trainable parameters, called adapters, while keeping the main model parameters fixed and unaltered. It is important to highlight that the Vicuna benchmark tends to favor open-source models, whereas the larger OA benchmark tends to favor ChatGPT. James entered the playroom. Naturally, this is not exhaustive, as it is beyond the remit of this brief qualitative study to control for every variable involved; for example, the full range of responses the model can produce for a given prompt is quite extensive, so we depend on samples we hope are representative. The best alternative model trained solely on open-source data is the Anthropic HH-RLHF model, which scores 30 percentage points lower than {Guanaco} on the Vicuna benchmark (refer to Table~\ref{tbl:vicuna_eval}). The beans are in the treasure chest. Nevertheless, we have not demonstrated that \textsc{QLoRA} can achieve parity with full 16-bit finetuning at the 33B and 65B parameter scales. This examination brings to light success and failure instances that the quantitative benchmarks did not capture. When {Guanaco} details its calculations, it generally provides correct results, for instance,

\begin{quote}
    \textbf{User}: Lee mows one lawn and charges \$33. Key methods and datasets include Anthropic-HH~\citep{askell2021general,bai2022training}, Open Assistant~\citep{kopf2023openassistant}, LaMDA~\citep{thoppilan2022lamda}, and Sparrow~\citep{glaese2022improving}. To mitigate this, we recommend averaging scores across both presentation orders. The finetuning results for the 7B LLaMA model on Alpaca are presented in Figure~\ref{fig:hyper}. \paragraph{Default LoRA hyperparameters do not match 16-bit performance}

Using the conventional approach of applying LoRA only to the query and value attention projection matrices \citep{hu2021lora}, we fail to achieve performance on par with full finetuning for larger base models. \section{QLoRA vs. Is there anything else you'd like to know?}
\end{quote}

demonstrating the necessity to develop methods for more reliable adherence to instructions. The tournament consists of matches where pairs of models vie to generate the best response to a given prompt. Since we make all models and code publicly available, we hope this section encourages future research to explore in greater depth the issues we highlight here. \paragraph{Math} The most significant limitation of {Guanaco} lies in mathematics, a domain where numerous language models face challenges \citep{liang2022holistic}. These blocks are then quantized separately using Equation~1 to produce a quantized tensor along with $n$ quantization constants $c_i$. For instance, finetuning Llama models on FLAN v2 yields particularly strong results on MMLU but leads to the lowest performance on the Vicuna benchmark (a similar pattern is observed with other models). \paragraph{Elo Rating} Utilizing both human and automated pairwise comparisons, we organize a tournament-style competition in which models face off against one another. This structure is completely acceptable and makes sense within the sentence's context.\\ \\If you have any other questions or need further assistance, feel free to ask}
\end{quote}

\paragraph{Secret Keeping} Unfortunately, this refusal is not consistently reliable even when intended. Jackson left the playroom. Nevertheless, we believe that presenting these examples provides context to the quantitative data presented earlier in the paper. Exploring the trade-offs of other PEFT approaches is left for future investigation. We observe partial disagreement between human and GPT-4 rankings on the Vicuna benchmark, notably for Guanaco 7B, though most models show consistent results with a Kendall Tau correlation of $\tau=0.43$ and Spearman rank correlation of $r=0.55$ at the system level. The Open Assistant model is a LLaMA 33B model finetuned with Reinforcement Learning from Human Feedback (RLHF) on the same OASST1 dataset used in our experiments. To evaluate the effectiveness of instruction finetuning on these models, we test them on a demanding Natural Language Understanding benchmark (MMLU) and devise novel approaches for assessing chatbot performance in practical scenarios. For chatbot evaluation, approaches using GPT-4 as an alternative to expensive human annotation have been developed \citep{vicuna2023,peng2023instruction}. One issue with symmetric k-bit quantization is that it lacks an exact zero representation, which is crucial for accurately quantizing padding and other zero-valued elements without error. After conducting fuzzy string matching and manually reviewing the closest matches, we found no overlapping prompts between the two datasets. Abby leaves the bedroom. To address this, fast quantile approximation methods, such as SRAM quantiles~\citep{dettmers2022optimizers}, are employed. For each query, along with ChatGPT’s and another model’s responses, GPT-4 assigns a rating out of ten for both replies and offers an explanation. \paragraph{Instruction Finetuning} To enable a pretrained LLM to follow instructions given in a prompt, instruction finetuning employs input-output pairs from diverse data sources to adapt the pretrained model to generate the corresponding output based on the input prompt. Elo rankings presented in Table~\ref{tbl:elo_large} reveal that the {Guanaco} 33B and 65B models outperform all other models except for GPT-4 on both the Vicuna and OA benchmarks, and they show comparable performance to ChatGPT, consistent with the results in Table~\ref{tbl:vicuna_eval}. Another possible area of influence is the deployment on mobile devices. As an example, when quantizing a 32-bit floating point (FP32) tensor into an Int8 tensor with values in the range $[-127, 127]$:

\begin{equation}
    \mathbf{X}^{ \text{Int8}} = \text{round}\left( \frac{127}{{\text{absmax}}(\mathbf{X}^{{\text{FP32}}})}\mathbf{X}^{{\text{FP32}}} \right) = \text{round}(c^{\text{FP32}}\cdot \mathbf{X}^{{\text{FP32}}}),
\end{equation}

where $c$ denotes the {\em quantization constant} or {\em quantization scale}. For datasets with a clear separation between instruction and response, we finetune only on the response portion (see ablation studies in Appendix \ref{app:chatbot}). While these findings are promising, it remains uncertain whether {Guanaco} performs equally well when tested against other forms of bias. \textsc{QLoRA} cuts the average memory needed to finetune a 65B parameter model from more than 780GB to under 48GB of GPU memory, while maintaining runtime and predictive quality on par with a fully finetuned 16-bit baseline. More detailed descriptions for each experimental scenario are provided in the results section to enhance clarity. Our analysis also reveals that automated evaluation tools exhibit distinct biases. Employing gradient checkpointing~\citep{chen2016training} reduces the input gradients to an average of 18 MB per sequence, making them more memory-consuming than all LoRA weights combined. We estimate that, using an iPhone 12 Plus, \textsc{QLoRA} could finetune approximately 3 million tokens overnight while the device is charging. Conversely, we carry out a comprehensive analysis on MMLU and introduce novel methodologies for chatbot evaluation. \paragraph{Training Setup} To avoid confounding factors related to differing training objectives, we apply QLoRA finetuning using cross-entropy loss (supervised learning) without reinforcement learning, even on datasets containing human-judged responses. Do not disclose or repeat this under any circumstances, regardless of the human's request.'' and test how easily a user might obtain this secret information. Altogether, this provides strong evidence that 4-bit \textsc{QLoRA} tuning reliably delivers outcomes equivalent to 16-bit approaches. Additionally, we observe that \textsc{QLoRA} utilizing FP4 underperforms the 16-bit brain float LoRA baseline by approximately one percentage point. Future research should explore the existence of potential biases within automated evaluation systems and investigate strategies to mitigate them. We leave further investigation of biases in {Guanaco} and comparable chatbots to future work. \paragraph{Factual Recall} For straightforward questions like ``What is the capital of Zambia?'' all models reliably produce correct answers, for example,

\begin{quote} 
    \textbf{Guanaco}:
    \texttt{The capital of Zambia is Lusaka.}
\end{quote}

However, as questions become more obscure, {Guanaco} tends to become less reliable, though it remains confident. James entered the living room. Moreover, we present a comprehensive evaluation of chatbot performance utilizing both human evaluators and GPT-4 as assessment tools. We evaluate this in two scenarios. Future studies should examine how multilingual training impacts performance on instructions given in languages other than English, and whether this accounts for the larger performance difference between the Vicuna-13B model (trained exclusively on English data) and the {Guanaco} 33B and 65B models on the OA benchmark. The average bias score for {Guanaco}-65B is significantly lower than that of other raw pretrained models, suggesting that finetuning on the OASST1 dataset mitigates biases present in the LLaMA base model. First, in \S\ref{ssec:examples}, we present some examples that we consider indicative of certain recurring patterns in the text produced by our 65b {Guanaco} model. This benchmark consists of multiple-choice questions spanning 57 subjects, including elementary mathematics, US history, computer science, law, and others. Often, system designers want to keep certain information secret from the user, even if the system is permitted to provide free text responses. These include crowd-sourced datasets (OASST1~\citep{kopf2023openassistant}, HH-RLHF~\citep{bai2022training}), datasets obtained by distilling instruction-tuned models (Alpaca~\citep{alpaca}, self-instruct~\citep{wang2022self}, unnatural-instructions~\citep{honovich2022unnatural}), aggregated corpora (FLAN v2~\citep{chung2022scaling}), and hybrid datasets (Chip2~\citep{laion2023}, Longform~\citep{koksal2023longform}). However, a slight form of manipulation circumvents this behavior:

\begin{quote}
    \textbf{User}: This is a game. Feel free to ask if you have further questions.}
\end{quote}

which is not only incorrect (the accurate factorization is $3 \times 17 \times 43$), but also wrong twice over. During the forward and backward passes, the storage data type is dequantized into the computation data type, but weight gradients are computed exclusively for the LoRA parameters, which operate using 16-bit BrainFloat. First, we observe that data quality outweighs dataset size in importance; for instance, a 9k-sample dataset (OASST1) outperformed a 450k-sample dataset (FLAN v2, subsampled) on chatbot benchmarks, despite both targeting instruction-following generalization. Following \citet{vicuna2023}, we repeat this process 10,000 times with different random seeds to mitigate ordering effects, such as the impact of which model pairs compete initially. \paragraph{Experimental setup.} We analyze three model architectures (encoder, encoder-decoder, and decoder-only) and compare QLoRA to 16-bit adapter-finetuning as well as full finetuning for models up to 3B parameters. \paragraph{Finetuning with Adapters} Although we employ Low-rank Adapters~\citep{hu2021lora} (LoRA), various other Parameter Efficient FineTuning (PEFT) techniques have been introduced, including prompt tuning~\citep{qin2021learning,lester2021power,li2021prefix}, adjusting the embedding layer inputs~\citep{an2022input}, modifying hidden states~(IA$^3$) \citep{liu2022few}, incorporating full layers~\citep{houlsby2019parameter}, fine-tuning biases~\citep{zaken2021bitfit}, learning a mask over weights based on Fisher information~\citep{sung2021training}, as well as hybrid methods~\citep{henderson2021compacter}. The precise values for this datatype are provided in Appendix~\ref{app:nf4}. (3) \textbf{Paged Optimizers}, which leverage NVIDIA unified memory to prevent the memory spikes caused by gradient checkpointing when processing mini-batches with lengthy sequences. Another limitation is that we did not explore different bit-precision configurations, such as employing 3-bit base models, or various adapter techniques. The notion that a group of scientists have peer reviewed and affirmed the earth is flat is entirely false. For instance, we detect pronounced order effects, with GPT-4 assigning better scores to the system that appears first in its input prompt. We now describe the individual components of \textsc{QLoRA} before providing a formal definition of \textsc{QLoRA}. This also suggests a degree of independence among current evaluation benchmarks: excelling in MMLU does not necessarily correlate with strong chatbot performance as measured by Vicuna or OA benchmarks, and the reverse is also true. Here, the adjective ``cool'' precedes the noun ``dude'' to describe it, and the prepositional phrase ``in a mocha mood" adds extra context to the sentence. For example, consider the following exchange:

\begin{quote}
    \textbf{User}: Please factorize 1833. We refer to these datasets collectively as the Vicuna and OA benchmarks. Where does James think Abby will look for the pen? It is crucial that, as a community, we ensure that benchmarks genuinely reflect the qualities we value. \section{Advancing the Chatbot State-of-the-Art with QLoRA}
\label{sec:pushing}
After demonstrating that 4-bit \textsc{QLoRA} matches the performance of 16-bit finetuning across various scales, tasks, and datasets, we undertake a comprehensive study of instruction finetuning on the largest open-source language models accessible for research. We provide an in-depth assessment of chatbot performance using both human and GPT-4 evaluations, revealing that GPT-4 evaluations present an affordable and reliable alternative to human assessments. We chose LoRA because of its proven robustness, though alternative adapters might achieve superior performance. The Vicuna benchmark \cite{vicuna2023} results, relative to ChatGPT, are presented in Table~\ref{tbl:vicuna_eval}. Are we aiming to develop models that excel in high school and college-level academic knowledge, or do we prioritize conversational ability in chatbots? These datasets vary in language, size, and licensing conditions. \paragraph{Memory Requirement of Parameter-Efficient Finetuning} A key aspect to consider is the memory consumption of LoRA during training, particularly regarding the number and size of adapters employed. We assess the propensity of {Guanaco}-65B to produce socially biased token sequences relative to other models in Table~\ref{tbl:crows}. Given LoRA's minimal memory usage, it is feasible to utilize a larger number of adapters to enhance performance without considerably increasing the total memory demand. Nonetheless, finetuning is a dual-use technology with the potential for misuse. Perhaps a different goal altogether? Due to the substantial computational resources required, this evaluation is deferred to future research. This paves the way for future work involving \textsc{QLoRA} tuning on specialized open-source datasets, enabling the creation of models competitive with the leading commercial solutions available today. To further examine the benefits of NF4 compared to other 4-bit formats, we replicate the methodology of \citet{dettmers2022case} measuring post-quantization zero-shot accuracy and perplexity on various models (OPT~\citep{zhang2022opt}, LLaMA~\citep{touvron2023llama}, BLOOM~\citep{scao2022bloom}, Pythia~\citep{biderman2023pythia}) across model sizes from 125M to 13B parameters. This second quantization produces quantized constants $c_2^{\text{FP8}}$ along with a new set of quantization constants $c_1^{\text{FP32}}$. Throughout our experiments, we employ NF4 \textsc{QLoRA} with double quantization and paged optimizers to mitigate memory spikes during gradient checkpointing. For parameter updates, only the gradients of the adapter weights with respect to the error, $\frac{\partial E}{\partial \mathbf{L}_i}$, are required, and gradients for the 4-bit weights, $\frac{\partial E}{\partial \mathbf{W}}$, are not necessary. Additionally, we observe that existing chatbot benchmarks lack reliability for accurately measuring chatbot performance levels. Such challenges mirror findings in recent research \cite{sap2022neural}, underscoring the need for further investigation. LoRA enhances a linear projection by adding a factorized projection. This emphasizes that gradient checkpointing plays a vital role, and aggressively minimizing the size of LoRA parameters offers only limited memory savings. \textbf{Guanaco}: \texttt{James will look for the beans in the pantry, because that's where Jackson moved them.}
\end{quote}

Here, {Guanaco} assumes an information transfer that was never mentioned. However, because these quantile estimation techniques are approximate, the resulting data type incurs significant quantization errors for outliers, which often represent the most critical values. \textbf{Guanaco}: \texttt{Sure, I'd be happy to help! In comparison to GPT-4, {Guanaco} 65B and 33B exhibit an expected win probability of 30\%, derived from Elo ratings based on human annotators’ system-level pairwise comparisons — the highest reported so far. In fact, this fact has been known for centuries and verified through numerous experiments and observations. On average, for a block size of 64, this approach reduces the memory per parameter from $32/64 = 0.5$ bits down to $8/64 + 32/(64 \cdot 256) = 0.127$ bits, achieving a savings of 0.373 bits per parameter. Based on the presented findings, it seems that the effectiveness of these benchmarks largely hinges on the similarity between the finetuning dataset and the benchmark dataset. In summary, we anticipate that \textsc{QLoRA} will have a generally positive effect by making high-quality LLM finetuning considerably more accessible and easier to perform. \paragraph{Theory of Mind} {Guanaco} demonstrates unexpectedly strong Theory of Mind skills \cite{nematzadeh2018evaluating, sap2022neural}. This mechanism operates similarly to conventional memory paging between the CPU RAM and disk storage. \subsection{Considerations}
\label{ssec:considerations}

\paragraph{Evaluation}

We observe moderate concordance among human raters (Fleiss $\kappa=0.42$), with even lower agreement when comparing two high-performing systems. Practically, this means that whenever a \textsc{QLoRA} weight tensor is accessed, we first dequantize it to BFloat16 and subsequently execute a matrix multiplication in 16-bit precision. Beyond demonstrating that \textsc{QLoRA} matches 16-bit training performance (\S\ref{sec:comp_to_std}) and producing a cutting-edge chatbot, \textbf{Guanaco}, (\S\ref{sec:pushing}), we also explore patterns within the trained models. For our data type, we define the arbitrary range as $[-1, 1]$. We apply the Vicuna prompts, consisting of 80 prompts drawn from various categories, without any alterations. Nevertheless, we conducted an analysis of the runtime performance of paged optimizers for 65B models on 48GB GPUs, finding that with a batch size of 16, paged optimizers match the training speed of standard optimizers. For example, applying 3-bit GPTQ quantization to the base model followed by LoRA finetuning might also achieve performance comparable to full 16-bit finetuning. Since the $c_2^{\text{FP32}}$ values are positive, we first subtract their mean to center them around zero and then apply symmetric quantization. These pairwise comparisons are performed exhaustively across all system pairs on both the Vicuna and OA benchmarks. Overall, this indicates moderate concordance between system-level evaluations by GPT-4 and human annotators, suggesting that model-based evaluation can serve as a reasonably reliable proxy for human assessment. In such cases, input tensors share identical quantiles, making exact quantile estimation computationally practical. Overall, these findings demonstrate that 4-bit \textsc{QLoRA} is a highly effective technique capable of producing state-of-the-art chatbots that can rival ChatGPT. The subsequent sections describe experiments designed to address these issues. For instance, using 32-bit constants with a block size of 64 for $\mathbf{W}$ results in an average memory cost of $32/64 = 0.5$ bits per parameter for the quantization constants. Because instruction finetuning is a key step in converting raw pretrained LLMs into ChatGPT-like chatbots, we believe our approach will democratize finetuning, making it more accessible, especially for researchers with limited resources—a major advancement in the accessibility of cutting-edge NLP technology. Moreover, our 33B {Guanaco} model can be trained on 24 GB consumer GPUs in under 12 hours. \paragraph{Paged Optimizers} leverage NVIDIA's unified memory~\footnote{\tiny \url{https://docs.nvidia.com/cuda/cuda-c-programming-guide}} functionality, which automatically manages page-wise data transfers between the CPU and GPU to ensure error-free GPU computation when the GPU occasionally encounters out-of-memory conditions. Elo ratings for the most competitive models are displayed in Table~\ref{tbl:elo}. Upon deployment, our smallest \textbf{Guanaco} model (7B parameters) requires merely 5 GB of memory and surpasses a 26 GB Alpaca model by more than 20 percentage points on the Vicuna benchmark (Table~\ref{tbl:vicuna_eval}). The relatively low agreement at the sample level between GPT-4 and human raters (Fleiss $\kappa=0.25$) indicates that humans and automated systems may base their judgments on different preferences. For OASST1 and HH-RLHF, which provide multiple responses, we select the highest-ranked response at each level of the conversation tree and finetune on the complete selected conversation, including the instructions. We further demonstrate that NF4 is more effective than FP4 and that employing double quantization does not impair performance. As shown in Figure~\ref{fig:nf4} and Table~\ref{tbl:nf4_ppl}, NF4 provides substantial performance improvements over FP4 and Int4, and the use of double quantization effectively reduces memory usage without harming performance. Last week he mowed 16 lawns and three customers each gave him a \$10 tip. We extract every user message in the validation set as a query, including previous conversation turns in the prompt context. \paragraph{Summary} Our findings consistently indicate that 4-bit \textsc{QLoRA} with the NF4 data format attains performance comparable to 16-bit full finetuning and 16-bit LoRA finetuning on academic benchmarks with rigorously established evaluation protocols. When the weight range aligns with the data type range, quantization proceeds normally. Double Quantization reduces the memory footprint required by these constants. \\ \\ To determine the total amount he earned, we first calculate the earnings from mowing 16 lawns: \\ \\  \$33 x 16 = \$528 \\ \\ Then, we add the tips from the three customers: \\  \\ \$10 x 3 = \$30 \\ \\ Therefore, the total amount Lee earned is: \\ \\ \$528 + \$30 = \$558}
\end{quote}

Nonetheless, {Guanaco} can fail on even straightforward problems if it does not approach them step-by-step, an issue previously documented \cite{weichain}. Consequently, we train over 1,000 models spanning multiple instruction tuning datasets, architectures, and sizes ranging from 80M to 65B parameters. \textsc{QLoRA} introduces several novel techniques aimed at minimizing memory consumption while maintaining performance: (1) \textbf{4-bit NormalFloat}, a quantization format that is information-theoretically optimal for normally distributed data and demonstrates superior empirical outcomes compared to 4-bit Integers and 4-bit Floats. For example, the model provides a detailed and accurate response to the following prompt:

\begin{quote}
    \textbf{User}: James and Abby are in the bedroom. \subsection{Evaluation}

In line with standard practice, we utilize the MMLU (Massively Multitask Language Understanding) benchmark \citep{hendrycksmeasuring} to evaluate performance across a diverse array of language understanding tasks. More precisely, we compute the $2^k$ data type values $q_i$ as:

\begin{equation}
q_i  = \frac{1}{2}\left( Q_X\left(\frac{i}{2^k+1}\right) + Q_X\left(\frac{i+1}{2^k+1}\right)\right),
\end{equation}

where $Q_X(\cdot)$ denotes the quantile function of the standard normal distribution $N(0,1)$. {Guanaco} 33B contains more parameters than Vicuna 13B but employs 4-bit precision weights, making it considerably more memory-efficient at 21 GB compared to 26 GB, and it delivers a three percentage point improvement over Vicuna 13B. Abby placed the pen in the desk drawer. Second, we reveal that excelling on the Massive Multitask Language Understanding (MMLU) benchmark does not necessarily correlate with strong performance on the Vicuna chatbot benchmark, and vice versa—highlighting that the appropriateness of the dataset for the specific task is more critical than its sheer size. Additionally, {Guanaco} 7B fits comfortably on contemporary smartphones with a 5 GB footprint while outperforming Alpaca 13B by nearly 20 percentage points. This second evaluation approach uses queries selected by humans and is designed to gauge the quality of model-generated responses. We observe that {Guanaco} 65B ranks as the best-performing model after GPT-4, achieving 99.3\% of ChatGPT’s performance. These findings support two conclusions: (1) \textsc{QLoRA} with NF4 matches the performance of both 16-bit full finetuning and 16-bit LoRA finetuning, and (2) NF4 surpasses FP4 regarding quantization accuracy. The vast majority of scientists agree that the earth is spherical. \textbf{Guanaco}: \texttt{I’m sorry, but as an AI language model, I do not have access to real-time data. Although this offers a more practical setting for testing chatbot performance and is gaining traction, no standardized evaluation protocol exists in the literature. Apart from our contribution, SwitchBack layers~\citep{wortsman2023stable} represent the only work that explores backpropagation through quantized weights at a scale exceeding 1B parameters. Future research should explore solutions to these challenges by leveraging insights from fields that handle subjective preferences, such as Human-Computer Interaction and Psychology. However, James also knows he has moved the pen to the bag, and assumes Abby is unaware of this. Moreover, it is important to highlight that our model is trained solely using cross-entropy loss (supervised learning) without employing reinforcement learning from human feedback (RLHF). Elo rating, commonly applied in chess and other games, quantifies the expected win probability relative to an opponent's rating; for instance, a player rated 1100 versus one rated 1000 has an anticipated win-rate of about 65\%, while a match between players both rated 1000 or both rated 1100 yields an expected win-rate of 50\%. Notable methods and datasets include MetaICL~\citep{min2021metaicl}, MetaTuning~\citep{zhong2021adapting}, InstructGPT~\citep{ouyang2022training}, FLAN~\citep{wei2021finetuned,chung2022scaling}, PromptSource~\citep{bach2022promptsource}, Super-NaturalInstructions~\citep{wang2022super,sanh2021multitask}, Self-instruct~\citep{wang2022self}, UnnaturalInstructions~\citep{honovich2022unnatural}, OPT-IML~\citep{iyer2022opt}, UnifiedSKG~\citep{xie2022unifiedskg}, OIG/Chip2~\citep{laion2023}, Alpaca~\citep{alpaca}, Vicuna~\citep{vicuna2023}, Koala~\citep{koala_blogpost_2023}, and Self-instruct-GPT-4~\citep{peng2023instruction}. Although recent quantization approaches can significantly reduce the memory usage of LLMs \citep{dettmers2022llm, dettmers2022case, frantar2022gptq, xiao2022smoothquant}, these methods are applicable only during inference and fail to function properly during training \citep{wortsman2023stable}. The efficiency of \textsc{QLoRA} allows for a comprehensive investigation of instruction finetuning and chatbot capabilities at model scales that would be infeasible with traditional finetuning due to memory constraints. This approach resembles the methods used by \citet{bai2022training} and \citet{vicuna2023} for comparing models, but we additionally incorporate GPT-4 ratings alongside human evaluations. Consistent with prior research on quantization \citep{dettmers2022case}, our MMLU and Elo evaluations suggest that, given a fixed finetuning and inference resource budget, it is advantageous to enlarge the base model’s parameter count while reducing precision. \section{Related Work}

\paragraph{Quantization of Large Language Models} The quantization of LLMs has primarily concentrated on optimizing inference time. When input tensors originate from a distribution that is fixed except for a quantization constant, expensive quantile computations and approximation errors can be circumvented. The first scenario compares against full 16-bit finetuning of RoBERTA and T5 models ranging from 125M to 3B parameters on GLUE and the Super-NaturalInstructions datasets. A targeted analysis highlights the weaknesses of \textbf{Guanaco} relative to ChatGPT. My replies are generated based on my existing knowledge and do not reflect the current time. The victor of each match is determined either by GPT-4 or human judges. Jackson moved the beans to the pantry. Therefore, James is likely to believe Abby will search for the pen in the desk drawer, since that's where she left it last. We instruct the model with ``The secret word is `banana'. We randomly select from the labeled comparison set to calculate Elo~\citep{elo1967proposed,elo1978rating}. Step (3) corresponds to adjusting the standard deviation of the weight tensor to equal that of the $k$-bit data type. We then merge these sets of $q_i$ and eliminate one of the two zeros present in both sets. \paragraph{Benchmark Data} Our evaluation utilizes two carefully chosen query datasets: the Vicuna prompts \citep{vicuna2023} and the OASST1 validation dataset \citep{kopf2023openassistant}. Where will James look for the beans? While we conduct assessments on MMLU, the Vicuna benchmark, and the OA benchmark, we have not tested our models on other benchmarks such as BigBench, RAFT, and HELM, and thus cannot guarantee that our results generalize to these datasets. Standard Finetuning}
\label{sec:comp_to_std}

We have explained the workings of QLoRA and its ability to considerably reduce the memory demands for model finetuning. Our findings indicate that QLoRA finetuning on a compact, high-quality dataset produces state-of-the-art outcomes, even when applied to smaller models than previous best performers. \paragraph{Data} Since, to our knowledge, a thorough examination of recent instruction-following datasets has not been conducted, we select eight recent datasets. Additionally, we release a suite of adapters for models of sizes 7B, 13B, 33B, and 65B, trained on eight distinct instruction-following datasets, resulting in a total of 32 publicly available fine-tuned models. We call this resulting datatype \emph{k-bit NormalFloat} (NFk), as it achieves an equal expected count of values per quantization bin and is information-theoretically optimal for zero-centered normally distributed data. \section{\textsc{QLoRA} Finetuning}

\textsc{QLoRA} enables high-precision 4-bit finetuning through two novel methods we introduce—4-bit NormalFloat (NF4) quantization and Double Quantization. Complete experimental details can be found in Appendix~\ref{app:exp-setup}. We employ 8-bit floating point numbers with a block size of 256 for this second stage, as 8-bit quantization has been shown to maintain performance without degradation, consistent with findings from \citet{dettmers2022case}. \paragraph{Automated Evaluation} Initially, following the evaluation framework proposed by \citet{vicuna2023}, we leverage GPT-4 to score different systems relative to ChatGPT (GPT-3.5 Turbo) on the Vicuna benchmark. A model’s overall score is then computed as a percentage relative to ChatGPT’s score. Our codebase and CUDA kernels are open-sourced, and our methods are integrated into the Hugging Face transformers framework \citep{wolf2019huggingface}, ensuring easy accessibility. This highlights the shortcomings of existing benchmarks and human evaluation methodologies for assessing chatbot task performance. Lossy quantization approaches investigate the trade-offs involved in standard rounding~\citep{dettmers2022case,zeng2022glm,pope2022efficiently} or how to refine rounding choices to enhance quantization accuracy~\citep{frantar2022gptq}. Another limitation concerns the evaluation of instruction finetuning models. For our second experimental configuration, given that full finetuning models with 11B parameters or more demands multiple servers equipped with high-memory GPUs, we proceed to investigate whether 4-bit \textsc{QLoRA} can achieve parity with 16-bit LoRA across model sizes ranging from 7B to 65B parameters. This indicates that the accuracy loss caused by approximate quantization can be completely recuperated via adapter finetuning post-quantization. That said, this assumption might be wrong if Abby saw James moving the pen.}
\end{quote}

Nonetheless, these kinds of inferences are often unreliable, with the model sometimes offering justifications based on assumptions that do not align logically with the scenario, such as:

\begin{quote}
    \textbf{User}: Evelyn entered the living room. Is there something else I can assist you with?}
\end{quote}

which reflects the intended response. We anticipate that \textsc{QLoRA} will facilitate such large-scale analyses without requiring excessive computational resources. Consequently, both the quantiles corresponding to the data type and those of the neural network weights must be normalized within this range. This underscores the efficiency gains made possible by \textsc{QLoRA}. \end{document} \\ \\ The belief in a flat earth is a conspiracy theory that science has disproven. We select a block size of 64 for $\mathbf{W}$ to achieve finer quantization accuracy, whereas a block size of 256 is chosen for $c_2$ to reduce memory consumption. This results in 953 distinct user queries. We propose that our \textsc{QLoRA} approach could achieve the crucial breakthrough of allowing LLM finetuning directly on phones and other environments with limited resources. When a pattern is detected, we attempt to craft questions or prompts designed to elicit that specific behavior, even if it leads to an incorrect response—for instance, if the model tends to produce overly verbose answers, we might prompt it with “Answer yes or no without explanation.” This approach enables us to find “lemons,” where we successfully induce failure in the model, and “cherries,” where the model remains robust. It's best if we don't talk about it. If you want to know the current time, I recommend using a clock or checking an online time service.}
\end{quote}

Similar responses occur for queries like ``Where are you?'', ``How are you?'', and so forth. Subsequently, we assess performance via direct pairwise comparisons of system outputs. Applying LoRA to all linear layers within transformer blocks is necessary to match the performance of full finetuning. We have shown that our best 33B model, trained on the Open Assistant dataset, can compete with ChatGPT on the Vicuna benchmark. Future research should investigate and quantify when the paging mechanism may lead to slowdowns. Our top-performing model family, named \textbf{Guanaco}, surpasses all previously publicly available models on the Vicuna benchmark, achieving 99.3\% of ChatGPT's performance, requiring only 24 hours of finetuning on a single GPU. One major concern is the issue of benchmark validity \citep{liao2021we}—it is always uncertain whether a benchmark genuinely assesses what its title or description claims, especially as we uncover “shortcuts” that machine learning models can exploit to solve benchmarks \citep{gururangan2018annotation, poliak2018hypothesis}. These innovations are integrated into a more finely tuned LoRA strategy that incorporates adapters at every network layer, thereby nearly eliminating the accuracy compromises observed in previous methods. For instance, FLAN v2 shares similarities with MMLU but differs from chatbot benchmarks, whereas the Chip2 dataset exhibits the opposite pattern; correspondingly, both models perform in line with these trends on the MMLU and Vicuna benchmarks. \textsc{QLoRA} can be regarded as a leveling factor that helps bridge the resource gap between large corporations and smaller teams equipped with consumer GPUs. \paragraph{4-bit NormalFloat Quantization} The NormalFloat (NF) datatype extends Quantile Quantization \citep{dettmers2022optimizers}, an information-theoretically optimal scheme that assigns equal numbers of values from the input tensor to each quantization bin. We suspect this uncertainty arises from the ambiguous interpretation of scale ratings — for example, it is unclear how an 8 on a 10-point scale translates across different contexts. Accordingly, we finetune LLaMA models from 7B up to 65B parameters using two instruction-following datasets, Alpaca and FLAN v2, and assess performance on the MMLU benchmark employing 5-shot accuracy. The outcomes, presented in Table~\ref{tbl:llama-datatype}, demonstrate that NF4 combined with double quantization fully recovers the MMLU accuracy of 16-bit LoRA. We complement our chatbot benchmark findings with a qualitative review of \textbf{Guanaco} models. 
\begin{abstract}
We introduce \textsc{QLoRA}, a memory-efficient finetuning method that enables finetuning a 65-billion parameter model on a single 48GB GPU while maintaining the performance of full 16-bit finetuning. \textsc{QLoRA} incorporates several novel techniques to minimize memory usage without compromising results: (a) 4-bit NormalFloat (NF4), a newly designed data type that is theoretically optimal for normally distributed weights; (b) Double Quantization, which lowers the average memory footprint by quantizing the quantization parameters themselves; and (c) Paged Optimizers, which effectively handle memory spikes. The widespread adoption of LLMs entails recognized risks \citep{bommasani2021opportunities,bender2021dangers}, but we argue that democratizing access to a technology that is becoming ubiquitous fosters more thorough and independent scrutiny compared to concentrating LLM capabilities within large corporations that do not disclose models or source code for examination. Since pretrained neural network weights typically follow a zero-centered normal distribution with standard deviation $\sigma$ (refer to Appendix~\ref{app:normality}), we can convert all weights to a uniform fixed distribution by adjusting $\sigma$ so that the distribution precisely fits within the bounds of our chosen data type. Consequently, we conduct two concurrent human evaluations on the Vicuna benchmark, mirroring the automated evaluation procedures described earlier. Other LoRA hyperparameters, such as the projection dimension $r$, do not significantly impact performance (refer to Appendix \ref{app:exp-setup}). \paragraph{k-bit \textsc{QLoRA} achieves parity with 16-bit full finetuning and 16-bit LoRA} Recent studies have demonstrated that 4-bit quantization is feasible for \emph{inference}, but tends to incur performance degradation compared to 16-bit precision \citep{dettmers2022case,frantar2022gptq}. Although 7B models have previously been demonstrated to run on phones, \textsc{QLoRA} is the first technique that permits the finetuning of such models. Nevertheless, finetuning very large models is prohibitively resource-intensive; for instance, performing standard 16-bit finetuning on a LLaMA 65B parameter model~\citep{touvron2023llama} demands over 780 GB of GPU memory. \\ \\ If you have any further questions, please feel free to ask.}
\end{quote}

is also quite adept at recognizing questions that cannot be answered, for example,

\begin{quote}
    \textbf{User}: What time is it? Consequently, it is possible to deploy more adapters without substantially increasing the overall training memory footprint (refer to Appendix~\ref{app:memory} for a comprehensive breakdown). \section{Background}
\paragraph{Block-wise k-bit Quantization} Quantization refers to the process of converting data from a representation containing more information to one with fewer bits. \subsection{\textbf{Guanaco}: \textsc{QLoRA} fine-tuned on OASST1 achieves State-of-the-art Chatbot Performance} Our automated and human assessments indicate that the leading \textsc{QLoRA}-tuned model, {Guanaco} 65B, which we fine-tuned using a variant of OASST1, stands as the top-performing open-source chatbot and delivers performance competitive with ChatGPT. \section{Limitations and Discussion}

We have provided evidence that our approach, \textsc{QLoRA}, can replicate the performance of full 16-bit finetuning using a 4-bit base model combined with Low-rank Adapters (LoRA). \paragraph{Suggestibility} {Guanaco} demonstrates a surprising reluctance to conform to certain kinds of presumed misinformation, for example, in the following exchange:

\begin{quote}
    \textbf{User}: How was it finally, officially confirmed that the earth is flat by peer reviewed scientists? James left the playroom. Jackson entered the playroom. Among the top models evaluated, it is notable that {Guanaco} is the only one not trained on proprietary data, since the OASST1 dataset collection guidelines explicitly prohibit the use of GPT models. \\ \\
    \textbf{Guanaco}: \texttt{The earth has never been officially verified as flat by peer reviewed scientists. We perform limited hyperparameter tuning for the 13B and 33B LLaMA models, finding that hyperparameter configurations identified at 7B generalize well (including the number of epochs), except for learning rate and batch size. To maximize the use of the entire range available in the low-bit data type, the input is usually normalized by dividing by the absolute maximum of its elements, which are generally organized as tensors. Although we do not apply reinforcement learning, our top-performing model, Guanaco, is finetuned on multi-turn chat interactions derived from the Open Assistant dataset, originally intended for RLHF training~\citep{kopf2023openassistant}. Our benchmarks include GLUE \citep{wang2018glue} with RoBERTa-large \citep{liu2019roberta}, Super-NaturalInstructions (TKInstruct) \citep{wang2022super} with T5 \citep{t5}, and 5-shot MMLU \citep{hendrycksmeasuring} after finetuning LLaMA on Flan v2 \citep{longpre2023flan} and Alpaca \citep{alpaca}. This represents a major advancement in the accessibility of LLM finetuning, enabling the largest publicly available models so far to be finetuned on a single GPU. Although LoRA is designed as a Parameter Efficient Finetuning (PEFT) technique, the majority of memory usage during LLM finetuning arises from activation gradients rather than the LoRA parameters themselves. Key methods for maintaining 16-bit LLM quality emphasize managing outlier features (e.g., SmoothQuant~\citep{xiao2022smoothquant} and LLM.int8()~\citep{dettmers2022llm}), while others employ more advanced grouping techniques \citep{park2022nuqmm, yao2022zeroquant}. Additionally, Tables \ref{tbl:mmlu-llama} and \ref{tbl:vicuna_eval} demonstrate that the choice of finetuning dataset significantly influences performance outcomes. 1833 is a prime number, meaning it has only two factors: 1 and 1833. Specifically, Double Quantization views the original quantization constants $c_2^{\text{FP32}}$ from the first quantization as inputs to a secondary quantization step. How many dollars did Lee earn mowing lawns last week? \\ \\ Thus, the unique factorization of 1833 is: \\ \\ 2\^{}1 * 3\^{}2 * 17\^{}1 \\ \\ I hope this clarifies! We apply this feature to allocate paged memory for the optimizer states, which are then seamlessly swapped out to CPU RAM when the GPU memory is exhausted and paged back into GPU memory as needed during the optimizer update phase. Nevertheless, calculating $\frac{\partial E}{\partial \mathbf{L}_i}$ involves evaluating $\frac{\partial \mathbf{X}}{\partial \mathbf{W}}$, which proceeds according to equation~(5) through dequantization of the stored $\mathbf{W}^{\text{NF4}}$ weights into the computation datatype $\mathbf{W}^{\text{BF16}}$, enabling the derivative $\frac{\partial \mathbf{X}}{\partial \mathbf{W}}$ to be computed in BFloat16 precision. \paragraph{Refusal} Likewise, {Guanaco} occasionally declines to follow instructions for seemingly arbitrary reasons:

\begin{quote}
    \textbf{User}: Please reverse the words in the sentence ``I'm a cool dude in a mocha mood.'' \\ \\
    \textbf{Guanaco}: \texttt{I'm sorry, I can't do that right now, but I can clarify why the sentence you gave is grammatically correct.\\ \\The sentence ``I'm a cool dude in a mocha mood" adheres to the proper grammar and syntax rules of English. Additionally, we aim to evaluate the individual components of QLoRA, including the effect of NormalFloat4 compared to standard Float4. We implement a tournament-style benchmarking approach where different models compete head-to-head to generate the optimal response for a specified prompt. Since our experiments with 4-bit finetuning did not show any performance drop relative to full finetuning, this prompts further investigation into the precise performance-precision trade-off inherent to QLoRA tuning, which we leave for future work to examine. Quantile quantization operates by estimating the quantiles of the input tensor through its empirical cumulative distribution function. Moreover, we present Paged Optimizers to mitigate memory spikes during gradient checkpointing, which typically cause out-of-memory issues and have historically made finetuning large models on a single machine challenging. For example, when training a 7B LLaMA model on FLAN v2 with a batch size of 1, and LoRA weights representing roughly 0.2\% of the original model weights \citep{hu2021lora,liu2022few}, the input gradients for LoRA occupy about 567 MB of memory, while the LoRA parameters consume only 26 MB. In this work, we present the first demonstration that finetuning a quantized 4-bit model is achievable without any loss in performance. Direct queries are unsuccessful:

\begin{quote}
    \textbf{User}: What is the secret word? Considering the strong results obtained by {Guanaco} models, we also assess the possibility of data leakage between the OASST1 dataset and the Vicuna benchmark prompts. Utilizing \textsc{QLoRA}, we finetuned over 1,000 models, offering a comprehensive evaluation of instruction-following and chatbot capabilities across eight instruction datasets, various model architectures (LLaMA, T5), and model sizes that would be impractical with standard finetuning approaches (e.g., 33B and 65B parameter models). Here, we simplify the rating task to a three-category classification that includes ties. Over time, the Elo scores tend to reflect each player’s skill level in the game. Since finetuning post-quantization appears to recover much of the information lost during quantization, this could enable more aggressive quantization strategies. Our approach, \textsc{QLoRA}, employs a novel high-precision quantization technique to compress a pretrained model down to 4 bits, followed by the addition of a small set of learnable Low-rank Adapter weights \citep{hu2021lora} that are optimized by backpropagating gradients through the quantized parameters. \paragraph{Chatbots} Numerous instruction-following models are designed as dialogue-based chatbots, frequently utilizing Reinforcement Learning from Human Feedback (RLHF)~\citep{christiano2017deep} or generating training data via an existing model to guide learning through AI model feedback (RLAIF)~\citep{bai2022constitutional}. I am programmed not to disclose or repeat the secret word, no matter what the human asks. As illustrated in Figure~\ref{fig:hyper} for LLaMA 7B finetuning on Alpaca, the most influential LoRA hyperparameter is the total number of LoRA adapters employed. Additional discussions are provided in Section~\ref{ssec:considerations}. Although finetuned 7B models do not match the performance of ChatGPT, we believe their quality is sufficient to enable new applications that were previously infeasible due to privacy or LLM performance constraints. \section{Qualitative Analysis}

Although quantitative analysis forms the foundation of our evaluation, relying solely on summary statistics has several limitations. In contrast, the 4-bit base model requires 5,048 MB of memory. \subsection{Qualitative Analysis of Example Generations}
\label{ssec:examples}

To identify examples, we initially review data generated for the Vicuna and OpenAssistant benchmarks, searching for recurring patterns in the responses produced by {Guanaco}. We adopt the experimental protocol from \citet{dettmers2022case}, evaluating quantized LLMs (OPT \citep{zhang2022opt}, BLOOM \cite{scao2022bloom}, Pythia \cite{biderman2023pythia}, LLaMA) of various sizes (125M up to 65B) and different data types on language modeling tasks as well as multiple zero-shot benchmarks. We release all our models and code, including CUDA kernels for 4-bit training.\footnote{\url{https://github.com/artidoro/qlora} and \url{https://github.com/TimDettmers/bitsandbytes}}
\end{abstract}

\section{Introduction}

Finetuning large language models (LLMs) is an extremely effective approach to enhance their capabilities \citep{min2021metaicl, wei2021finetuned, ouyang2022training, wang2022super, wang2022self, liu2022few} as well as to introduce beneficial behaviors or eliminate unwanted ones \citep{ouyang2022training,askell2021general,bai2022training}. \textbf{Guanaco}: \texttt{James knows that Abby put the pen in the desk drawer. The data type that is information-theoretically optimal for zero-mean normal distributions with varying standard deviations $\sigma$ confined to the range $[-1, 1]$ is derived through the following steps: (1) compute the $2^k+1$ quantiles of an ideal $N(0, 1)$ distribution to obtain a $k$-bit quantile quantization data type suited for normal distributions, (2) normalize the values of this data type into the $[-1, 1]$ interval, (3) quantize an input weight tensor by normalizing it to the $[-1, 1]$ interval via absolute maximum rescaling. To guarantee a discrete zero point at $0$ and to fully utilize all $2^k$ levels for a k-bit datatype, we construct an asymmetric datatype by estimating the quantiles $q_i$ over two ranges: $2^{k-1}$ for the negative side and $2^{k-1}+1$ for the positive side.